% ============================================================
% RESUMEN
% ============================================================

\begin{abstract}

Este informe consolida los diez ejercicios de procesamiento digital de
señales desarrollados en \texttt{THEORY/FIRST\_ROUND/WORK/} (carpetas
\texttt{WORK\_ONE} a \texttt{WORK\_TEN}), cada uno implementado como un
paquete Python autocontenido con su propia suite de pruebas. Se cubren:
síntesis y verificación espectral de tonos DTMF; reconstrucción de una
onda cuadrada mediante series de Fourier; análisis espectral (FFT, media,
desviación estándar y separabilidad estadística) de grabaciones de
animales, instrumentos y voces; convolución discreta con núcleos tipo
escalón y trapezoidal; análisis de una señal de modulación de amplitud a
dos frecuencias de muestreo; espectrogramas wavelet y un ejercicio de
correlación deslizante sobre audio; filtrado digital de promedio móvil a
múltiples frecuencias de corte; diseño manual de filtros IIR de segundo y
cuarto orden mediante la transformada bilineal, validado numéricamente
contra \texttt{scipy.signal}; modulación de amplitud de audio real a
50~kHz y 500~kHz de frecuencia de muestreo; y un banco de filtros
FIR/IIR combinado con la transformada de Hilbert, incluyendo un barrido
de orden de filtro. En todos los casos los resultados numéricos y las
figuras provienen de la ejecución real del código correspondiente, no de
datos sintéticos inventados; las limitaciones conocidas de cada ejercicio
(por ejemplo, la ausencia de una tercera grabación de voz, o la síntesis
por voz artificial en \texttt{WORK\_THREE}) se documentan explícitamente
en su respectiva sección.

\end{abstract}

% ============================================================
% PALABRAS CLAVE
% ============================================================

\begin{IEEEkeywords}

Procesamiento digital de señales, Transformada Rápida de Fourier, FFT,
convolución discreta, filtros digitales, filtros FIR, filtros IIR,
transformada de Hilbert, wavelets, modulación de amplitud.

\end{IEEEkeywords}
